\section{Affine Hilbert cubes and Hales--Jewett-free composition}\label{sec:cubes}

\subsection{Progression limits vs.\ Hales--Jewett in two-stage composition}\label{sec:hj-wall}

Suppose two linear class-uniform maps
$K_1, K_2 : \Mat_{d\times d}(F^b) \to \Mat_{d\times d}(F^b)$ with class partitions
$p_1 : \Z_d \to \mathrm{Fin}\,M_1$ and $p_2 : \Z_d \to \mathrm{Fin}\,M_2$ satisfy
$K_2 \circ K_1 = \tau$. By \cite[Section~6, Lemma~6.1 and Theorem~6.3]{SharedDecoders2026}, any
such factorization induces a common-message subspace $W \le A$ inside a $q$-channel cyclic module
$A \cong (\Z_d \to F^q)$ of dimension $\dim_F A \le b M_2 d$ ($q \le b M_2$) together with a linear
readout $L : A \to Z$ injective on $W$ such that $L(T_c w) = L(w)$ for all $w \in W$ and all $c$
in a single $p_1$-class $J = p_1^{-1}(m_1)$ (where $(T_s w)_u = w_{u+s}$, formalized as
\lean{cyclicChannelShift}). When $J$ contains a length-$\ell$ arithmetic progression
$\{c_0 + j s : 0 \le j < \ell\}$, iterating $T_s$ yields
$\ell \dim_F W \le \dim_F A + (\ell - 1)\dim_F\ker(T_s - \mathrm{id})$.

We record both the Hales--Jewett baseline and the Behrend/random progression barriers inside
\texttt{BoundedArityDecoders} to make the quantitative bottleneck precise.

\begin{proposition}[Hales--Jewett baseline vs.\ Behrend and random colorings]\label{prop:diag-hj-behrend}
\formal{DiagnosisAndBarriersProofs.lean}{diagnosis_RamseyHalesJewettTwoStage,diagnosis_BehrendAndRandomAPBound}
\formal{DiagnosisAndBridgesConsequences.lean}{diagnosis5_CombinatorialInvVdWTendstoAndRankBound,diagnosis6_ModularBehrendAndRandomZModProgression}
\formal{DiagnosisAndBarriers.lean}{Diagnosis_RamseyHalesJewettTwoStageSpec,Diagnosis_BehrendAndRandomAPBoundSpec}
\begin{enumerate}[label=(\alph*)]
\item \emph{Hales--Jewett progression baseline}
  (\cite{HalesJewett1963}; \cite[Section~6]{SharedDecoders2026}): for every fixed $M_1 \ge 1$, the
  largest monochromatic progression length $\ell(d, M_1)$ guaranteed in every $M_1$-coloring of
  $\Z_d$ tends to $\infty$ as $d \to \infty$, ruling out exact two-stage linear factorizations at
  fixed $M_1, M_2$ for sufficiently large $d$.
\item \emph{Behrend and random progression limits} (\cite{Behrend1946}): there exist $3$-AP-free
  subsets of $\Z_d$ of size $d / 2^{O(\sqrt{\log d})}$---hence $M_1$-colorings with
  $M_1 = 2^{O(\sqrt{\log d})}$ containing \emph{no} non-constant $3$-term progression modulo
  $d$---and $2$-colorings of $\Z_d$ containing no monochromatic progression of length
  $\Omega(\log d)$.
\end{enumerate}
\end{proposition}
\begin{outline}
Part~(a) derives divergence of the inverse van der Waerden threshold from Mathlib's Hales--Jewett
theorem and plugs $\ell > M_2$ into the common-message dimension bound. Part~(b) embeds base-$B$
Euclidean spheres of $\{0,\dots,\lfloor(B-1)/2\rfloor\}^n$ into $\Z_d$ without wraparound carries
(\cite{Behrend1946}) and bounds monochromatic progressions in random $2$-colorings by a union bound.
\end{outline}

\subsection{The Affine \texorpdfstring{$k$}{k}-Cube Common-Message Growth Lemma}\label{sec:cube-lemma}

To dispense with arithmetic progressions over an \emph{arbitrary} field $F$, consider any $k$
linear endomorphisms $R_0, \dots, R_{k-1} \in \End_F(A)$ (not even assumed to commute) and define
for each $\eps \in \{0,1\}^k$ the ordered operator cube vertex (\lean{orderedCubeMap})
\[
  R^\eps \;:=\; R_{k-1}^{\eps_{k-1}} \cdots R_0^{\eps_0} \in \End_F(A).
\]

\begin{lemma}[Affine $k$-cube common-message growth]\label{lem:affine-cube}
\formal{AffineCubeCompositionProofs.lean}{lemma_AffineCubeCommonMessageGrowth}
\formal{AffineCubeComposition.lean}{Lemma_AffineCubeCommonMessageGrowthSpec}
Let $F$ be any field, $A$ a finite-dimensional $F$-vector space, $Z$ an $F$-vector space,
$R_0, \dots, R_{k-1} \in \End_F(A)$, $W \le A$ a subspace, and $L : A \to Z$ an $F$-linear map
injective on $W$. If $L(R^\eps x) = L(x)$ for all $x \in W$ and all $\eps \in \{0,1\}^k$, then
\[
  (k + 1)\dim_F W \;\le\; \dim_F A \;+\; \sum_{t=0}^{k-1} \dim_F \ker(R_t - \mathrm{id}_A).
\]
\end{lemma}
\begin{outline}
Define the nested cube-orbit subspaces $S_0 := W$ and $S_{t+1} := S_t + R_t(S_t)$ for
$0 \le t < k$. By induction on $t$, every $u \in S_t$ is a linear combination of elements
$R^\eps x$ with $x \in W$ and $\eps_s = 0$ for $s \ge t$. Appending $\eps_t = 1$ versus
$\eps_t = 0$ in the cube invariance $L(R^\eps x) = L(x)$ shows that $P_t := R_t - \mathrm{id}_A$
maps $S_t$ into $S_{t+1} \cap \ker L$. Since $W \subseteq S_t \subseteq S_{t+1}$ and $L$ is
injective on $W$, we have $W \cap P_t(S_t) \subseteq W \cap \ker L = \{0\}$, so
\[
  \dim_F S_{t+1} \;\ge\; \dim_F W + \dim_F P_t(S_t)
  \;\ge\; \dim_F W + \dim_F S_t - \dim_F \ker(R_t - \mathrm{id}_A).
\]
Summing from $t = 0$ to $k - 1$ gives $(k+1)\dim_F W \le \dim_F S_k + \sum_{t=0}^{k-1}
\dim_F\ker(R_t - \mathrm{id}_A) \le \dim_F A + \sum_{t=0}^{k-1} \dim_F\ker(R_t - \mathrm{id}_A)$.
\end{outline}

\subsection{Elementary Cauchy--Schwarz cube extraction and Theorem 2A}\label{sec:thm2a}

Unlike arithmetic progressions, affine $k$-cubes are extracted in a single subset $S \subseteq \Z_d$
(even when $S$ is $3$-AP-free!) by elementary bucket pigeonholing and Cauchy--Schwarz, with explicit
step bounds $1 \le s_t \le \lfloor d / (2 k^2)\rfloor$ that guarantee $\dim_F\ker(T_{s_t} - \mathrm{id})$
stays small.

\begin{lemma}[Cauchy--Schwarz affine $k$-cube extraction and dense-class bound]\label{lem:cs-cube}
\formal{AffineCubeCompositionProofs.lean}{cauchySchwarzAffineCubeExtraction}
\formal{AffineCubeConsequences.lean}{singleSubset_CauchySchwarzAffineCubeExtraction,denseSubset_AffineCubeCommonMessageBound}
\formal{AffineCubeComposition.lean}{CauchySchwarzAffineCubeExtractionSpec}
Let $d, M_{\mathrm{eff}}, k \ge 1$ satisfy $\max\bigl(8 M_{\mathrm{eff}} k^2,\;
(2 M_{\mathrm{eff}})^{2^k}\bigr) < d$.
\begin{enumerate}[label=(\alph*)]
\item \emph{Single-subset form}: for every subset $J \subseteq \Z_d$ of density $\abs{J} \ge
  d / M_{\mathrm{eff}}$ (regardless of whether $J$ is $3$-AP-free), there exist $c_0 \in J$ and
  steps $s_0, \dots, s_{k-1} \in \{1, \dots, \lfloor d/(2k^2)\rfloor\}$ such that
  $c_0 + \sum_{t : \eps_t = 1} s_t \in J$ for all $\eps \in \{0,1\}^k$. Over any field $F$, any
  subspace $W \le (\Z_d \to F^q)$ with a linear readout injective on $W$ and shift-invariant across
  $J$ satisfies $(k + 1)\dim_F W \le q d + q k \lfloor d / (2 k^2)\rfloor$, ruling out exact
  two-stage linear factorizations whenever a single stage-1 class over an arbitrary label type
  $\kappa_1$ has size $\ge d / M_{\mathrm{eff}}$ at $k = M_2$
  (\lean{denseSubset_AffineCubeCommonMessageBound}).
\item \emph{Coloring form}: setting $M_{\mathrm{eff}} = M_1$ and $k = M_2$, every $M_1$-coloring
  $p_1 : \Z_d \to \mathrm{Fin}\,M_1$ admits a monochromatic affine $M_2$-cube with steps in
  $\{1, \dots, \lfloor d/(2 M_2^2)\rfloor\}$.
\end{enumerate}
\end{lemma}
\begin{outline}
Lift $J$ to $S \subseteq \{0,\dots,d-1\}$. Partitioning $\{0,\dots,d-1\}$ into $Q = 2k^2$ intervals
of length $\le \lfloor d/Q\rfloor + 1$ and applying Cauchy--Schwarz to the bucket sizes shows that
some step $s \in \{1, \dots, \lfloor d/Q\rfloor\}$ satisfies
$\abs{S}(\abs{S} - Q) \le 2d\,\abs{\{x \in S : x + s \in S\}}$. Iterating $k$ times shrinks the
effective density denominator according to $D_{t+1} = 2 D_t^2$ ($D_0 = M_{\mathrm{eff}}$, so
$D_k = \frac{1}{2}(2 M_{\mathrm{eff}})^{2^k}$); the threshold $d > \max(8 M_{\mathrm{eff}} k^2,
(2M_{\mathrm{eff}})^{2^k})$ keeps every iterate nonempty and prevents wraparound modulo $d$ because
$\sum_{t=0}^{k-1} s_t \le k \lfloor d/(2k^2)\rfloor < d$. Feeding the extracted cube into
\Cref{lem:affine-cube} with $\dim_F\ker(T_{s_t} - \mathrm{id}) \le q s_t \le q\lfloor d/(2k^2)\rfloor$
yields the common-message dimension bound.
\end{outline}

\begin{theorem}[Theorem 2A: Hales--Jewett-free two-stage exact factorization cutoff]\label{thm:2a}
\formal{AffineCubeCompositionProofs.lean}{theorem2A_TwoStageRamseyFreeCutoff}
\formal{AffineCubeConsequences.lean}{theorem2A_GrowingClassesNonFactorization}
\formal{AffineCubeComposition.lean}{Theorem2A_TwoStageRamseyFreeCutoffSpec}
Let $F$ be \emph{any} field, let $d, b, M_1, M_2 \ge 1$, and let
$p_1 : \Z_d \to \mathrm{Fin}\,M_1$ and $p_2 : \Z_d \to \mathrm{Fin}\,M_2$. If two $F$-linear
class-uniform maps $K_1, K_2$ on $\Mat_{d\times d}(F^b)$ satisfy $K_2 \circ K_1 = \tau$, then
\[
  d \;\le\; \max\!\Bigl(8 M_1 M_2^2,\; (2 M_1)^{2^{M_2}}\Bigr).
\]
Consequently, for any growing class schedules $(M_1(d), M_2(d))$ eventually satisfying
\[
  \max\!\Bigl(8 M_1(d) M_2(d)^2,\; \bigl(2 M_1(d)\bigr)^{2^{M_2(d)}}\Bigr) \;<\; d
\]
---in particular for fixed $M_2$
and $M_1(d) = 2^{O(\sqrt{\log d})}$ on $3$-AP-free Behrend partitions, or for fixed $M_1$ and
$M_2(d) \le (1 - \eps)\log_2\log_2 d$---no exact two-stage linear factorization $K_2 \circ K_1 =
\tau$ exists over any field $F$.
\end{theorem}
\begin{outline}
Assume for contradiction that $\max(8M_1 M_2^2, (2M_1)^{2^{M_2}}) < d$ and $K_2 \circ K_1 = \tau$.
By \Cref{lem:cs-cube}(b), some $p_1$-class contains an affine $M_2$-cube with steps
$1 \le s_t \le \lfloor d/(2M_2^2)\rfloor$. On the $q$-channel cyclic common-message module
$A = (\Z_d \to F^q)$ ($q \le b M_2$, $\dim_F A = qd \le b M_2 d$), the column-message reduction
(\cite[Lemma~6.1 and Theorem~6.3]{SharedDecoders2026}) produces a subspace $W \le A$ of dimension
$\dim_F W = bd$
while $\dim_F\ker(T_{s_t} - \mathrm{id}) = q\gcd(s_t, d) \le q s_t \le b M_2 \lfloor d/(2M_2^2)\rfloor$.
Applying \Cref{lem:affine-cube} with $k = M_2$ gives
\[
  (M_2 + 1) bd \;\le\; b M_2 d + \sum_{t=0}^{M_2 - 1} b M_2 \lfloor d/(2 M_2^2)\rfloor
  \;\le\; b M_2 d + \frac{1}{2} bd,
\]
contradicting $bd > 0$.
\end{outline}

\subsection{Sharpness at \texorpdfstring{$d = 13$}{d = 13} and same-interval coded factorizations at \texorpdfstring{$d = 4$}{d = 4}}\label{sec:cube-controls}

Two finite-dimensional controls establish that both halves of \Cref{lem:affine-cube} and
\Cref{thm:2a} are genuinely sharp.

\begin{example}[Sharpness of $(k+1)$ vs.\ $2^k$ on an AP-free $2$-cube at $d = 13$]\label{ex:d13-sharpness}
\formal{AffineCubeCompositionProofs.lean}{control_D13TwoCubeAPFreeSharpness}
\formal{AffineCubeConsequences.lean}{control_D13_GeneralFieldAndProgressionBridge}
\formal{AffineCubeComposition.lean}{Control_D13TwoCubeAPFreeSharpnessSpec}
Let $d = 13$, $q = 2$ ($\dim_F A = 26$), and consider the $2$-cube
$J_{13} = \{0, 1, 4, 5\} = \{0,1\} + \{0,4\} \subseteq \Z_{13}$.
\begin{enumerate}[label=(\alph*)]
\item $J_{13}$ is $3$-AP-free modulo $13$ (\lean{ThreeAPFreeMod}), so every non-constant
  progression in $J_{13}$ has length $\ell \le 2$, where the progression formula
  $\lfloor(26 + 2(\ell - 1))/\ell\rfloor$ gives only $\dim_F W \le 14$ (attained on $\{0,1\}$ over
  every field).
\item Over \emph{every} field $F$, \Cref{lem:affine-cube} with $k = 2$ and
  $\dim_F\ker(T_s - \mathrm{id}) = 2$ ($s \in \{1,4\}$) gives $3\dim_F W \le 26 + 2 + 2 = 30$, so
  $\dim_F W \le 10 < 13$.
\item Conversely, over $\R$---and more generally over every field $F$ with $(13 : F) \ne 0$---there
  exists an explicit $(W, L)$ invariant on $J_{13}$ of dimension $\dim_F W = 9$. Since
  $4 \cdot 9 = 36 > 26 + 3 \cdot 2 = 32$, the naive $2^k$-multiplier inequality
  $2^k \dim_F W \le \dim_F A + (2^k - 1)q$ ($\dim_F W \le 8$) is \emph{false}.
\end{enumerate}
\end{example}
\begin{outline}
Exhaustive residue check verifies $J_{13}$ is $3$-AP-free, and cyclic decomposition of $\Z_{13}$
gives $\dim_F\ker(T_s - \mathrm{id}) = 2$ for $s \not\equiv 0 \pmod{13}$ over every field. For the
$D = 9$ witness, the constant frequency ($2$ dimensions) plus $7$ non-constant Fourier modes whose
$2\times 2$ channel amplitudes satisfy a $3$-constraint $J_{13}$-alignment system inside a
$4$-dimensional mode space contribute $2 + 7 = 9$ independent dimensions (realized over $\R$ via
real/imaginary cosine-sine pairs and over general $(13 : F) \ne 0$ in
\lean{control_D13_GeneralFieldAndProgressionBridge}).
\end{outline}

\begin{example}[Same-interval coded factorizations at $d = 4 = M_1 M_2$]\label{ex:d4-coded}
\formal{AffineCubeCompositionProofs.lean}{control_D4SameIntervalCodedFactorization}
\formal{AffineCubeConsequences.lean}{control_D4_AllWidthsAndBinaryFields}
\formal{AffineCubeComposition.lean}{Control_D4SameIntervalCodedFactorizationSpec}
Let $d = 4$ and let both stages use the \emph{identical} two-interval partition
$p_1 = p_2 = [0, 0, 1, 1]$ (\lean{halfIntervalPartition4}, $M_1 = M_2 = 2$).
\begin{enumerate}[label=(\alph*)]
\item \emph{Uncoded impossibility}: for every lane width $b \ge 1$ and finite alphabet
  $\abs{\Sigma} \ge 2$, neither entry-copy rules nor bit-lane-copy rules $K_1, K_2$ with
  $p_1 = p_2 = [0,0,1,1]$ can satisfy $K_2 \circ K_1 = \tau$.
\item \emph{General-field coded factorization ($(20 : F) \ne 0$)}: over $\R$, and over every field
  $F$ with $(20 : F) \ne 0$, for every lane width $b \ge 0$ there exist explicit $F$-linear
  class-uniform maps $K_1, K_2$ on $\Mat_{4\times 4}(F^b)$ with $p_1 = p_2 = [0,0,1,1]$ satisfying
  exact transposition $K_2 \circ K_1 = \tau$.
\item \emph{Characteristic-$2$ and $(5 : F) \ne 0$ coded factorization across all widths (including
  $b = 1$ over $\F_2$)}: in characteristic~$2$ ($(2 : F) = 0$, including $\F_2$), a distinct pair of
  scalar $0/1$ coefficient tables achieves $K_2 \circ K_1 = \tau$ lane-by-lane for \emph{every}
  lane width $b \ge 0$ (including scalar $b = 1$ over $\F_2$). Combining $(2 : F) = 0$ and
  $(20 : F) \ne 0$ yields exact linear coded factorizations $K_2 \circ K_1 = \tau$ at all widths $b$
  over every field $F$ with $(5 : F) \ne 0$ (\lean{control_D4_AllWidthsAndBinaryFields}).
\end{enumerate}
\end{example}
\begin{outline}
In~(a), projecting onto a single lane reduces bit-lane copy to scalar entry copy over $\Sigma$, and
pigeonholing the $4$ entries of a row through $2$-element interval classes forces a routing
collision whenever linear coding across matrix positions is forbidden. In~(b) and~(c), linear
combinations across entries within each $4\times 4$ scalar lane disambiguate the two-element
intervals $\{0,1\}$ and $\{2,3\}$: integer tables scaled by $20^{-1}$ satisfy the exact transpose
identity whenever $(20 : F) \ne 0$, while separate binary $0/1$ tables satisfy the exact transpose
identity modulo~$2$ whenever $(2 : F) = 0$; any field with $(5 : F) \ne 0$ has either $(2 : F) = 0$
or $(20 : F) \ne 0$.
\end{outline}
